\documentclass[11pt, a4paper]{article}

% bfw-style.tex — gemeinsame Style-Definitionen für alle Handouts/Aufgabenhefte
% Voraussetzung: Kompilation mit xelatex (wegen fontspec)

\usepackage[ngerman]{babel}
\usepackage{geometry}
\geometry{a4paper, top=2.2cm, bottom=2.2cm, left=2.3cm, right=2.3cm}

\usepackage{fontspec}
% Fallback-Kette: Source Sans Pro -> Calibri -> Segoe UI -> Latin Modern Sans
% (Latin Modern ist immer mit MiKTeX vorhanden)
\IfFontExistsTF{Source Sans Pro}{%
  \setmainfont{Source Sans Pro}[Ligatures=TeX]%
}{\IfFontExistsTF{Calibri}{%
  \setmainfont{Calibri}[Ligatures=TeX]%
}{\IfFontExistsTF{Segoe UI}{%
  \setmainfont{Segoe UI}[Ligatures=TeX]%
}{%
  \setmainfont{Latin Modern Sans}[Ligatures=TeX]%
}}}
\IfFontExistsTF{Source Code Pro}{%
  \setmonofont{Source Code Pro}[Scale=0.92]%
}{\IfFontExistsTF{Consolas}{%
  \setmonofont{Consolas}[Scale=0.92]%
}{%
  \setmonofont{Latin Modern Mono}[Scale=0.92]%
}}

\usepackage{xcolor}
% Farbpalette: hell, freundlich, modern
\definecolor{bfwPrimary}{HTML}{0EA5A4}    % Petrol/Türkis - Primär
\definecolor{bfwPrimaryDark}{HTML}{0F766E} % dunkler Petrol für Headings
\definecolor{bfwAccent}{HTML}{F59E0B}     % Amber - Akzent
\definecolor{bfwSuccess}{HTML}{10B981}    % Grün - Erfolg / Lösung
\definecolor{bfwWarn}{HTML}{F97316}       % Orange - Achtung
\definecolor{bfwInk}{HTML}{0F172A}        % fast Schwarz
\definecolor{bfwInkSoft}{HTML}{475569}    % gedämpftes Grau
\definecolor{bfwBgSoft}{HTML}{F8FAFC}     % off-white
\definecolor{bfwBgTeal}{HTML}{ECFEFF}     % helles Türkis
\definecolor{bfwBgAmber}{HTML}{FEF3C7}    % helles Gelb
\definecolor{bfwBgRose}{HTML}{FFE4E6}     % helles Rosé für Eselsbrücken

\usepackage[colorlinks=true, linkcolor=bfwPrimaryDark, urlcolor=bfwPrimaryDark]{hyperref}
\usepackage{enumitem}
\setlist[itemize]{leftmargin=*, itemsep=2pt, topsep=2pt}
\setlist[enumerate]{leftmargin=*, itemsep=2pt, topsep=2pt}

\usepackage[most]{tcolorbox}
\usepackage{booktabs}
\usepackage{tikz}
\usetikzlibrary{decorations.pathmorphing, positioning, shapes.geometric, arrows.meta}

% Merkkästchen — wichtige Definition / Kernpunkt
\newtcolorbox{merksatz}[1][]{%
  enhanced, breakable,
  colback=bfwBgTeal,
  colframe=bfwPrimary,
  boxrule=0.6pt,
  arc=4pt,
  left=10pt, right=10pt, top=8pt, bottom=8pt,
  fonttitle=\bfseries\color{bfwPrimaryDark},
  title={\faLightbulb\ Merksatz},
  #1
}

% Eselsbrücke — Mnemoniks
\newtcolorbox{eselsbruecke}[1][]{%
  enhanced, breakable,
  colback=bfwBgRose,
  colframe=bfwAccent,
  boxrule=0.6pt,
  arc=4pt,
  left=10pt, right=10pt, top=8pt, bottom=8pt,
  fonttitle=\bfseries\color{bfwWarn},
  title={Eselsbrücke},
  #1
}

% Beispiel-Box
\newtcolorbox{beispiel}[1][]{%
  enhanced, breakable,
  colback=bfwBgSoft,
  colframe=bfwInkSoft,
  boxrule=0.4pt,
  arc=3pt,
  left=10pt, right=10pt, top=8pt, bottom=8pt,
  fonttitle=\bfseries\color{bfwInk},
  title={Beispiel},
  #1
}

% Lösungs-Box (für Aufgabenhefte)
\newtcolorbox{loesung}[1][]{%
  enhanced, breakable,
  colback=bfwBgAmber!40,
  colframe=bfwSuccess,
  boxrule=0.6pt,
  arc=3pt,
  left=10pt, right=10pt, top=8pt, bottom=8pt,
  fonttitle=\bfseries\color{bfwSuccess!70!black},
  title={Lösung},
  #1
}

% Glossar-Eintrag
\newcommand{\glossareintrag}[2]{%
  \par\noindent\textbf{\color{bfwPrimaryDark}#1} \enspace #2\par\smallskip
}

% Section-Stil — etwas farbiger als Standard
\usepackage{titlesec}
\titleformat{\section}
  {\Large\bfseries\color{bfwPrimaryDark}}
  {\thesection}{0.6em}{}
\titleformat{\subsection}
  {\large\bfseries\color{bfwPrimary}}
  {\thesubsection}{0.5em}{}
\titleformat{\subsubsection}
  {\normalsize\bfseries\color{bfwInk}}
  {\thesubsubsection}{0.4em}{}

% TikZ-Stil "skizzenhaft" — etwas weicher als Rough.js, aber stabil
\tikzset{
  roughbox/.style = {
    draw=bfwPrimaryDark, thick, fill=bfwBgTeal,
    rounded corners=4pt, inner sep=6pt
  },
  roughaccent/.style = {
    draw=bfwAccent, thick, fill=bfwBgAmber,
    rounded corners=4pt, inner sep=6pt
  },
  roughline/.style = {
    draw=bfwInkSoft, thick, -{Stealth[length=2.5mm]}
  }
}

% Bullet-Symbole (bewusst kein FontAwesome — vermeidet Font-Installations-Probleme)
\newcommand{\faLightbulb}{$\bullet$}
\newcommand{\faBookmark}{$\bullet$}

% Header/Footer
\usepackage{fancyhdr}
\pagestyle{fancy}
\fancyhf{}
\renewcommand{\headrulewidth}{0pt}
\fancyfoot[L]{\small\color{bfwInkSoft}\thepage}
\fancyfoot[R]{\small\color{bfwInkSoft} BFW M-064 Beschaffung im Büromanagement}


\title{\vspace*{-1.5cm}\Huge\bfseries\color{bfwPrimaryDark}Tag~9: Vertiefung und Reflektion\\[0.4em]
  \large\color{bfwInkSoft}\textnormal{Wiederholung aller Themen + IHK-Cheatsheet (Tag 1--8)}}
\author{\textsf{Beschaffung im Büromanagement}\\
\textsf{\small\copyright\ Can Siebert\,\textperiodcentered\,2026}}
\date{22.05.2026}

\begin{document}
\maketitle
\thispagestyle{empty}

\vspace{1em}
\begin{merksatz}[title={\bfseries Worum geht es heute?}]
Acht Tage Beschaffung im Büromanagement liegen hinter Ihnen. Heute bringen wir
alles zusammen --- die wichtigsten Konzepte, Formeln und Paragraphen kompakt,
ergänzt um typische Querverbindungen und IHK-Klausurmuster. Dieses Handout dient
als \emph{Cheatsheet} für die Prüfungsvorbereitung.
\end{merksatz}

\tableofcontents
\vspace{1em}
\hrule
\vspace{1em}

\section{Beschaffungsprozess --- die Kette}

\begin{merksatz}[title={Die Schritte}]
\textbf{Bedarf} $\to$ \textbf{Bezugsquelle} $\to$ \textbf{Anfrage} $\to$ \textbf{Angebot} $\to$ \textbf{Bestellung} $\to$ \textbf{Lieferung} $\to$ \textbf{Wareneingang} $\to$ \textbf{Rechnung} $\to$ \textbf{Zahlung}
\end{merksatz}

Bedarfsarten: Roh-, Hilfs-, Betriebsstoffe. \emph{6~W der Bestellung}: Was, Wer,
Wann, Wo, Wie viel, Welche Qualität. \emph{ABC-Analyse}: Konzentration auf wenige
A-Güter mit hohem Wertanteil.

\section{Bezugskalkulation --- die wichtigste Rechnung im Einkauf}

\begin{merksatz}[title={Bezugskalkulation}]
Listenpreis $-$ Liefererrabatt $=$ Zieleinkaufspreis $-$ Liefererskonto $=$
\textbf{Bareinkaufspreis} $+$ Bezugskosten $=$ \textbf{Bezugspreis}
\end{merksatz}

\begin{beispiel}
Listenpreis 1\,000~€, 10\,\% Rabatt $\to$ 900~€; 2\,\% Skonto $\to$ 882~€;
Bezugskosten 18~€ $\to$ Bezugspreis $\mathbf{900~€}$.
\end{beispiel}

\section{Vertragsschluss und AGB}

\begin{itemize}[itemsep=2pt]
  \item Antrag (§ 145 BGB) + Annahme (§ 147 BGB) $=$ Vertrag.
  \item \emph{Essentialia negotii}: Kaufgegenstand, Kaufpreis, Vertragsparteien.
  \item AGB §§ 305 ff.\ BGB --- Einbeziehung, überraschende Klauseln (§ 305c), Inhaltskontrolle (§ 307).
  \item Verbraucher (§ 13) vs.\ Unternehmer (§ 14) --- entscheidet über anwendbares Recht.
\end{itemize}

\section{Wareneingang und § 377 HGB}

\begin{merksatz}[title={Wareneingang in 4 Schritten}]
\textbf{Annahme} $\to$ \textbf{Äußere Kontrolle} (vor dem Fahrer) $\to$ \textbf{Innere Kontrolle} (nach Abfahrt) $\to$ \textbf{Erfassung}
\end{merksatz}

\textbf{§ 377 HGB Rügepflicht beim Handelskauf}: Käufer muss unverzüglich
untersuchen und Mängel anzeigen. Versäumnis $=$ Genehmigung $=$ Verlust aller
Mängelrechte.

\section{Lagerkennzahlen --- Klausurklassiker Teil 2}

\begin{merksatz}[title={Die vier Formeln}]
\begin{enumerate}[itemsep=0pt]
  \item \textbf{Ø Lagerbestand} $=$ (AB $+$ EB) $\div$ 2
  \item \textbf{LUH} $=$ Jahresverbrauch $\div$ Ø Lagerbestand
  \item \textbf{Ø Lagerdauer} $=$ 360 $\div$ LUH
  \item \textbf{Lagerzinsen} $=$ Ø Lagerwert $\times$ Lagerzinssatz
\end{enumerate}
\end{merksatz}

Hohe LUH $=$ wenig Kapitalbindung $=$ wirtschaftlich gut.

\section{Kaufvertragsstörungen}

\subsection{Schuldnerverzug --- § 286 BGB}

\textbf{Voraussetzungen}: Fälligkeit + Mahnung + Vertretenmüssen.
\textbf{Mahnung entbehrlich} bei kalendermäßiger Bestimmung oder § 286 III BGB
(30 Tage nach Fälligkeit + Rechnung bei Geldschulden).

\textbf{Verzugszinsen} (§ 288 BGB):
$\text{Zinsen} = \dfrac{\text{Forderung} \times \text{Zinssatz} \times \text{Tage}}{360}$.
B2B: Basiszins $+$ 9\,\%-Punkte; B2C: $+$ 5\,\%-Punkte.

\subsection{Annahmeverzug --- §§ 293 ff.\ BGB}

\textbf{Voraussetzungen}: ordnungsgemäßes Angebot + Nichtannahme.
\textbf{Folgen}: Haftungsmilderung (§ 300), Mehraufwendungen (§ 304), Gefahrübergang.

\subsection{Schlechtleistung --- § 437 BGB}

\begin{merksatz}[title={Käuferrechte-Treppe}]
\textbf{1.\ Nacherfüllung} (§ 439, Vorrang!) $\to$ \textbf{2.\ Rücktritt} (§ 323)
\textbar{} \textbf{3.\ Minderung} (§ 441) \textbar{} \textbf{4.\ Schadensersatz} (§§ 280, 281)
\end{merksatz}

\textbf{Minderung}: Kaufpreis $\times \dfrac{\text{Wert mangelhaft}}{\text{Wert mangelfrei}}$.

\section{Eingangsrechnungen --- Eigentum, Gefahr, Prüfung}

\subsection{Eigentumsvorbehalt --- § 449 BGB}

Vier Formen: \emph{einfach} / \emph{erweitert} (Kontokorrent) / \emph{verlängert}
(Wiederverkauf $+$ Vorausabtretung) / \emph{weitergeleitet}.

\subsection{Gefahrübergang --- §§ 446, 447 BGB}

\textbf{§ 446}: Übergabe $\to$ Gefahr beim Käufer (Grundregel).
\textbf{§ 447}: Versendungskauf B2B $\to$ Gefahr ab Übergabe an Spediteur.
\textbf{§ 475 II}: Verbrauchsgüterkauf $\to$ Gefahr erst beim Verbraucher.

\subsection{Pflichtangaben einer Rechnung --- § 14 UStG (8 Punkte)}

\begin{enumerate}[itemsep=0pt]
  \item Adressen Lieferant + Empfänger
  \item Steuernummer / USt-IdNr.
  \item Ausstellungsdatum
  \item Fortlaufende Rechnungsnummer
  \item Menge + handelsübliche Bezeichnung
  \item Liefer-/Leistungsdatum
  \item Nettoentgelt nach Steuersätzen
  \item Steuersatz und Steuerbetrag
\end{enumerate}

\subsection{Rechnungsprüfung}

Drei Stufen: \emph{formell} (Pflichtangaben) $\to$ \emph{sachlich} (Bestellung,
Lieferung) $\to$ \emph{rechnerisch} (Beträge). Drei-Punkt-Vergleich:
Bestellung $\to$ Wareneingang/Lieferschein $\to$ Rechnung.

\section{Bezahlung --- SEPA, Karten, Skonto}

\subsection{SEPA}

Überweisung $|$ Basis-Lastschrift (B2C, 8 Wochen Erstattung) $|$ Firmen-Lastschrift
(B2B, kein Erstattungsrecht) $|$ Dauerauftrag.

\subsection{Skonto --- der IHK-Klassiker}

\begin{merksatz}[title={Effektivverzinsung der Skonto-Inanspruchnahme}]
$\text{Effektivverzinsung} = \dfrac{\text{Skontosatz} \times 360}{\text{Zahlungsfrist} - \text{Skontofrist}}$
\end{merksatz}

Beispiel ,,2/14 --- netto 30'': $\dfrac{2 \times 360}{30 - 14} = 45\,\%$ p.\,a.
Skonto fast immer ziehen.

\section{Querverbindungen --- so denkt die IHK}

\begin{itemize}[itemsep=2pt]
  \item Wareneingang (Tag 5) + Mängelrüge § 377 HGB (Tag 5) $\to$ Käuferrechte § 437 BGB (Tag 6) $\to$ Verzug § 286 BGB (Tag 6).
  \item Bestellung (Tag 3) $\to$ Lieferung (Tag 4) $\to$ Eigentumsvorbehalt § 449 BGB (Tag 7) sichert die Forderung des Lieferanten.
  \item Rechnung (Tag 7) $\to$ Pflichtangaben § 14 UStG $\to$ Vorsteuerabzug $\to$ Skonto-Entscheidung (Tag 8).
  \item Lagerkennzahlen (Tag 5) $\leftarrow$ Bezugspreis (Tag 2) liefert die Bezugswerte für Lagerwert/Zinsberechnung.
\end{itemize}

\section{Top-10 zum Auswendiglernen}

\begin{enumerate}[itemsep=0pt]
  \item Bezugskalkulation: Listenpreis $-$ Rabatt $-$ Skonto $+$ Bezugskosten $=$ Bezugspreis.
  \item § 377 HGB --- unverzügliche Mängelrüge beim Handelskauf.
  \item LUH $=$ Jahresverbrauch $/$ Ø~Lagerbestand.
  \item Ø~Lagerdauer $=$ 360 $/$ LUH.
  \item § 286 BGB --- Verzug $=$ Fälligkeit $+$ Mahnung $+$ Vertretenmüssen.
  \item Verzugszinsen B2B: Basiszins $+$ 9~\%-Punkte; B2C $+$ 5~\%-Punkte.
  \item § 437 BGB --- Käuferrechte (Nacherfüllung vor Rücktritt/Minderung/Schadensersatz).
  \item § 447 BGB --- Versendungskauf B2B: Gefahr ab Spediteur.
  \item § 14 UStG --- 8 Pflichtangaben einer Rechnung.
  \item Skonto-Effektivzins $=$ Skontosatz $\times$ 360 $/$ (Zahlungsfrist $-$ Skontofrist).
\end{enumerate}

\newpage

\section*{Glossar (Cheatsheet)}
\addcontentsline{toc}{section}{Glossar}

\glossareintrag{Beschaffungsprozess}{9-stufige Kette von Bedarfsermittlung bis Zahlung.}

\glossareintrag{Bezugskalkulation}{Schema zur Ermittlung des Bezugspreises aus Listenpreis, Rabatten, Skonti, Bezugskosten.}

\glossareintrag{ABC-Analyse}{Konzentration der Beschaffung auf wenige A-Güter mit hohem Wertanteil.}

\glossareintrag{6 W der Bestellung}{Was, Wer, Wann, Wo, Wie viel, Welche Qualität.}

\glossareintrag{Anfrage}{allgemein (offene Konditionen) vs.\ bestimmt (konkrete Anfrage).}

\glossareintrag{Antrag (§ 145 BGB)}{verbindlich, ohne ,,freibleibend''-Klausel.}

\glossareintrag{AGB (§§ 305 ff.\ BGB)}{Allgemeine Geschäftsbedingungen --- Einbeziehung, Inhaltskontrolle.}

\glossareintrag{§ 377 HGB}{Rügeobliegenheit beim Handelskauf --- unverzüglich nach Lieferung.}

\glossareintrag{Mängelrüge}{schriftliche, eindeutige Anzeige eines festgestellten Mangels.}

\glossareintrag{Lagerumschlagshäufigkeit (LUH)}{Jahresverbrauch / Ø~Lagerbestand --- wie oft das Lager pro Jahr umgeschlagen wird.}

\glossareintrag{Ø Lagerdauer}{360 / LUH --- wie lange Ware im Schnitt im Lager liegt.}

\glossareintrag{Lagerzinsen}{Ø Lagerwert × Lagerzinssatz --- Opportunitätskosten der Kapitalbindung.}

\glossareintrag{Schuldnerverzug (§ 286 BGB)}{Fälligkeit + Mahnung + Vertretenmüssen.}

\glossareintrag{Verzugszinsen (§ 288 BGB)}{B2B: Basiszins + 9 \%-Punkte; B2C: + 5 \%-Punkte.}

\glossareintrag{Annahmeverzug (§§ 293 ff.\ BGB)}{Gläubiger nimmt ordnungsgemäß angebotene Leistung nicht an.}

\glossareintrag{§ 437 BGB}{Käuferrechte: Nacherfüllung, Rücktritt, Minderung, Schadensersatz.}

\glossareintrag{Nacherfüllung (§ 439 BGB)}{Nachbesserung oder Ersatzlieferung --- Vorrang vor allen anderen Rechten.}

\glossareintrag{Minderung (§ 441 BGB)}{Kaufpreis × (Wert mangelhaft / Wert mangelfrei).}

\glossareintrag{Eigentumsvorbehalt (§ 449 BGB)}{einfach / erweitert / verlängert / weitergeleitet.}

\glossareintrag{§ 932 BGB}{gutgläubiger Eigentumserwerb vom Nichtberechtigten --- nicht bei abhandengekommenen Sachen (§ 935).}

\glossareintrag{§ 446 BGB}{Gefahrübergang mit Übergabe (Grundregel).}

\glossareintrag{§ 447 BGB}{Gefahrübergang ab Übergabe an Spediteur (Versendungskauf B2B).}

\glossareintrag{§ 14 UStG}{Pflichtangaben einer Rechnung --- Voraussetzung für Vorsteuerabzug.}

\glossareintrag{Drei-Punkt-Vergleich}{Abgleich Bestellung -- Wareneingang/Lieferschein -- Rechnung.}

\glossareintrag{SEPA}{Single Euro Payments Area --- Überweisung, Lastschrift, Dauerauftrag.}

\glossareintrag{Skonto-Effektivverzinsung}{Skontosatz × 360 / (Zahlungsfrist − Skontofrist) --- meist sehr hoch.}

\end{document}
